
\section{定轴转动中的功能关系}

\begin{definition}[力矩的功]
	
刚体定轴运动中，外力做功仅包括垂直转轴平面内分量。

已知力矩 $M_i  = F_irisin\phi \qquad$  角位移为$ \theta \qquad$

则合外力做功为$$A = \sum A_i = \sum \int_{\theta_0}^{\theta} M_i d\theta = \int_{\theta_0}^{\theta}(\sum M_i)d\theta = \int_{\theta_0}^{\theta}Md\theta \qquad $$ 
式中 $M = \sum M_i$ 为刚体所受到的合外力矩。

\end{definition}

\begin{definition}[刚体转动惯量]
$J = \sum m_ir_i^2$
\end{definition}

\begin{definition}[刚体转动动能]
$E_k = \frac{1}{2} J \omega^2$
\end{definition}

\begin{definition}[刚体定轴转动动能定理]
合外力矩对刚体所做的功等于刚体动能增量
\end{definition}
\begin{supplement}
工程上的应用：转动惯量很大的飞轮储存转动动能
\end{supplement}

